\section*{Multiselection (author source pp.11:1--11:2)}
In the selection problem, also commonly called multiplexing, one is given a long string $x\in\bits^n$ and an index $i\in[n]$, and the goal is to output $x_i$. Selection is one of the basic operations in the word RAM model (and indeed costs unit time in that model). Implementing it by a Boolean circuit however is more involved and requires a circuit of size $\Theta(n)$.
In this work we investigate a natural generalization of selection, which we call multiselection; now rather than a single index, one is given indices $i_1,\ldots,i_q\in[n]$, and the goal is to compute $(x_{i_1},\ldots,x_{i_q})$. We denote this function by $\Sel^{n\to q}$. Multiselection is trivial in the word RAM model of computation (it costs $\Theta(q)$ operations with any standard instruction set). We focus on implementing multiselection by small Boolean circuits, where here and throughout this work all circuits have bounded fan-in. In this setting, we do not have matching upper and lower bounds. The best circuit size lower bound for multiselection is the input length $n+q\cdot\log n$. On the other hand, the straightforward construction (which selects each index separately) is much larger, with size $\Theta(n\cdot q)$. A slightly more sophisticated approach reduces to sorting, and results in a circuit of size $\Theta(n\log^2n)$ for any $q=O(n)$. Somewhat surprisingly, it appears that multiselection has not been systematically studied and in particular we are not aware of essentially any other upper bounds (see further discussion in Section 1.1).
